% !TeX root = ../../Book1.tex
\OptionalChapter{可整流集与可整流性}{b1:ch:19}% 本章摘要（不计入编号）
\begin{chapterabstract}
面积与余面积公式已经说明，Lipschitz 参数片可以承载微分和积分。本章转而从集合与测度本身出发，定义可数可整流性和纯不可整流性，证明局部有限面积集合的近似切空间与单位密度，再建立切向微分、面积公式及带权推前。最后讨论可整流／纯不可整流分解，并以明确的文献边界介绍密度判据与 Preiss 定理。
\end{chapterabstract}

% 本章目标（不计入编号）
\begin{learninggoals}
\begin{itemize}
\item\label{b1:item:ch19-goal-rectifiable}能区分参数像、可数可整流集、纯不可整流集及可整流测度，注意局部有限性和绝对连续条件。
\item\label{b1:item:ch19-goal-tangent}理解单片缩放、片间余量与近似切平面的关系，并能用图的微分计算切空间和法空间。
\item\label{b1:item:ch19-goal-area}能在参数片上定义与坐标选择无关的切向微分，正确使用切向面积及推前密度公式。
\item\label{b1:item:ch19-goal-density}理解密度为一的顺向结论、纯不可整流分解与 Preiss 逆向判据各自需要哪些工具。
\end{itemize}
\end{learninggoals}

本章为选读内容，但不是此前定理的补充假设：第十七、十八章的面积与余面积证明已经独立闭合。读者可以先掌握可整流性的定义和图上的切向积分，把最后一节的深层密度判据留作第二遍阅读。关于近似切平面的主线，先修是 Rademacher、面积公式、Lebesgue 点，以及第十四章对一般 Radon 测度的微分；后一个工具用来说明相邻参数片为何不会在典型点留下额外的主阶质量。

Simon 的《Lectures on Geometric Measure Theory》第11--12节\cite{Simon1983Lectures}是切空间与切向面积公式的主要参照。本书改用前章已经建立的 Lipschitz 分片，把参数独立性与片间测度比较展开。Mattila 的《Rectifiability: A Survey》\cite{Mattila2023Rectifiability}提供集合、测度与密度判据的系统对照；其测度定义有不要求相对于 Hausdorff 测度绝对连续的版本，本书会明确区分。Preiss 定理只给有出处的陈述与证明概要，其深层证明不作为本章已自行证完的内容。

% !TeX root = ../../Book1.tex
\section{可整流集}
\label{b1:sec:1901}

前两章总是先给一个映射，再计算它的像或纤维。本章改变出发点：只给一个集合，问它能否在忽略零测部分后由可数个 Lipschitz 参数块描述。这样既保留曲线和曲面的局部微分结构，又允许交叉、断裂和多张参数片。可整流性正是对这种测度意义下的几何结构所作的要求。

% 实读：Mattila2023Rectifiability 出版社OA正式全书§4.1，定义4.1--4.2及随后约定说明。
% 本书区分可数可整流与有限Hk测度；可整流测度另要求mu<<Hk，严于该书定义4.2。
% 基本闭合性质、紧片、四角Cantor纯不可整流例和RN表示均用本书前章工具完整展开。

\subsection{Lipschitz 像}
\label{b1:subsec:190101}

以下先固定整数 \(1\leq k\leq n\)。若 \(A\subseteq\mathbb R^k\)，且 \(f:A\to\mathbb R^n\) 为 Lipschitz 映射，则 \(f(A)\) 是一个由 \(k\) 个参数描述的像。不过，这个描述不保证参数唯一，也不保证微分满秩；一张参数片可以折回或塌缩到较低维集合。

定义可整流性时，可以把参数映射取为在整个 \(\mathbb R^k\) 上定义。因为逐坐标使用第十六章的延拓定理，任意 \(f:A\to\mathbb R^n\) 都可延为全空间的 Lipschitz 映射，至多增大一个仅依赖目标维数的常数。这里仅需存在一个有限常数，不要求向量值延拓保留最优常数。

全空间参数化还有一个可测性上的好处。若 \(f:\mathbb R^k\to\mathbb R^n\) Lipschitz，则
\[
 f(\mathbb R^k)=\bigcup_{j=1}^{\infty}f\bigl(\overline B(0,j)\bigr)
\]
是可数个紧集之并，因而 Borel 可测。每个紧像的 \(\mathcal H^k\) 测度有限，因为
\[
 \mathcal H^k\Bigl(f\bigl(\overline B(0,j)\bigr)\Bigr)
 \leq\operatorname{Lip}(f)^k\omega_kj^k.
\]
当 \(f\) 为常值时，左端为零，同一结论仍成立。因此可数参数化天然给出的是 \(\sigma\) 有限性；它不必给出全体像的有限面积，也不必给出每个环境球中的有限面积。

\subsection{可数 \texorpdfstring{\(m\)}{m}-可整流集}
\label{b1:subsec:190102}

\begin{definition}\label{b1:def:countably-rectifiable-set}
若集合 \(E\subseteq\mathbb R^n\) 是 \(\mathcal H^k\) 可测的，并且存在 Lipschitz 映射 \(f_j:\mathbb R^k\to\mathbb R^n\)，使
\[
 \mathcal H^k\left(E\setminus\bigcup_{j=1}^{\infty}f_j(\mathbb R^k)\right)=0,
\]
则称 \(E\) 为\term[keshu-kezhengliuji]{可数 \(k\) 维可整流集}[countably k-rectifiable set]。在本章不致混淆时简称 \(k\) 维可整流集。若还要求 \(\mathcal H^k(E)<\infty\)，会明确写出这一附加条件。
\end{definition}

定义允许一个 \(\mathcal H^k\) 零测余项，但不允许遗漏具有正 \(k\) 维测度的部分。零测集因而自动满足定义；这只是定义在零测情形下没有几何要求，并不意味着每个零测集都是一张平滑曲面。相反，一个集合即使不可数、拓扑复杂，也可能在所讨论维数下没有需要参数化的质量。

\begin{proposition}\label{b1:prop:rectifiable-basic-properties}
可数 \(k\) 维可整流集具有以下性质。
\begin{enumerate}
\item\label{b1:item:rectifiable-sigma-finite}其 \(\mathcal H^k\) 测度为 \(\sigma\) 有限；其中每个 \(\mathcal H^k\) 可测子集仍可整流。
\item\label{b1:item:rectifiable-countable-union}可数个这样的集合之并仍可整流；增删 \(\mathcal H^k\) 零测集不改变可整流性。
\item\label{b1:item:rectifiable-lipschitz-image}若 \(g:E\to\mathbb R^p\) 为 Lipschitz 映射，则 \(g(E)\) 是 \(\mathcal H^k\) 可测的可数 \(k\) 维可整流集，其中目标维数取 \(p\geq k\)。
\end{enumerate}
\end{proposition}
\begin{proof}
第一项的 \(\sigma\) 有限性来自上一小节的有界参数块及零测余项；对子集仍用原来的参数映射即可。第二项把各组可数参数映射并为一列，遗漏的可数个零测余项之并仍零测。

证明第三项时，先把 \(g\) 逐坐标延为 \(\mathbb R^n\) 上的 Lipschitz 映射 \(\widetilde g\)。于是 \(\widetilde g\circ f_j\) 给出所需参数映射，而遗漏部分的像由 Hausdorff 测度的 Lipschitz 估计仍为零测。

还须说明像的可测性。由第一项及 Hausdorff 测度的 Borel 包络、有限 Borel 测度的紧内正则性，可以取紧集 \(K_j\subseteq E\)，使 \(\mathcal H^k(E\setminus\bigcup_jK_j)=0\)。具体地，先将 \(E\) 分成有限测度块，在各块的有限 Borel 包络上使用正则性，再删去一个覆盖改值余项的 Borel 零测集即可取得块内紧逼近。各 \(g(K_j)\) 紧，余项的像零测，所以 \(g(E)\) 在完备 \(\mathcal H^k\) 下可测。
\end{proof}

特别地，\(C^1\) 图像和光滑水平集的满秩部分都可整流。前者在相对紧的小球上有界导数，因而局部 Lipschitz；用可数个这样的球覆盖参数域即可。后者由上一章的正规坐标命题局部化成图像，再利用欧氏空间的可数拓扑基选取可数覆盖。

上一章还给出了一个非光滑来源。若 \(f:\mathbb R^m\to\mathbb R^n\) Lipschitz 且 \(m>n\)，则其几乎每个水平集都可数 \((m-n)\) 维可整流。事实上，余面积证明中的满秩坐标片在固定目标值后由 \(\mathbb R^{m-n}\) 的子集作 Lipschitz 参数化，临界与不可微部分则在几乎每层上为零测。把各子集上的参数映射延到全空间，正好满足本定义。\(m=n\) 时，相应结论是几乎每个纤维至多可数，可按零维约定理解。

可整流性不包含局部有限面积的要求。例如可数条平行直线可以在平面中稠密，它们的并可数 \(1\) 维可整流，但每个非空开球内的总长度为无穷。下一节需要研究球内密度时，会另外要求 \(\mathcal H^k|_E\) 局部有限。不能只凭已经写下“可整流”三个字便自动使用这一条件。

\subsection{纯不可整流集}
\label{b1:subsec:190103}

\begin{definition}\label{b1:def:purely-unrectifiable-set}
若 \(P\subseteq\mathbb R^n\) 是 \(\mathcal H^k\) 可测的，且对每个 Lipschitz 映射 \(f:\mathbb R^k\to\mathbb R^n\) 都有
\[
 \mathcal H^k\bigl(P\cap f(\mathbb R^k)\bigr)=0,
\]
则称 \(P\) 为\term[chun-bukezhengliuji]{纯 \(k\) 维不可整流集}[purely k-unrectifiable set]。
\end{definition}

等价地，\(P\) 与每个可数 \(k\) 维可整流集的交都为 \(\mathcal H^k\) 零测集：只需把该可整流集的参数片逐个相交，再加上零测余项。这个条件比“\(P\) 不是可整流集”强得多，因为后者可能还包含一大块可整流部分。

\begin{proposition}\label{b1:prop:rectifiable-pure-zero}
若一个 \(\mathcal H^k\) 可测集既可数 \(k\) 维可整流又纯 \(k\) 维不可整流，则它为 \(\mathcal H^k\) 零测集。反过来，每个 \(\mathcal H^k\) 零测集都同时具有这两个性质。
\end{proposition}
\begin{proof}
把可整流集的参数覆盖代入纯不可整流的定义，每片所截到的测度都为零，加上零测余项便得到整个集合零测。反向结论由两个定义直接得到。
\end{proof}

下面给出一个有正长度却纯不可整流的紧集，说明这种分类不是仅在零测情形下成立。令 \(C\subseteq[0,1]\) 为保留左右各四分之一区间的 Cantor 集，并置 \(P=C\times C\)。

\begin{proposition}\label{b1:prop:four-corner-pure-unrectifiable}
上述四角 Cantor 集满足 \(0<\mathcal H^1(P)<\infty\)，并且纯 \(1\) 维不可整流。
\end{proposition}
\begin{proof}
第 \(j\) 步的 \(P\) 包含于 \(4^j\) 个边长 \(4^{-j}\) 的正方形中，直径覆盖的总代价为 \(\sqrt2\)。所以 \(\mathcal H^1(P)\leq\sqrt2\)。第十五章 Cantor 集构造中的等权概率测度 \(\nu\) 满足 \(\nu(I)\leq C_0|I|^{1/2}\)；这里对应 \(\lambda=1/4\)，其维数为 \(1/2\)。对乘积概率 \(\mu=\nu\times\nu\)，半径为 \(r\) 的球包含于两条长度为 \(2r\) 的区间的乘积内，因而
\[
 \mu\bigl(\overline B(x,r)\bigr)\leq2C_0^2r.
\]
由质量分布原理，\(\mathcal H^1(P)>0\)。

证明纯不可整流性。任取 Lipschitz 曲线 \(\gamma:\mathbb R\to\mathbb R^2\)，写 \(\gamma=(\gamma_1,\gamma_2)\)，并令 \(A=\gamma^{-1}(P)\)。这个集合闭，且 \(\gamma_i(A)\subseteq C\)。由于 \(m_1(C)=0\)，对标量映射 \(\gamma_i\) 使用一维面积公式，得到
\[
 \int_A|\gamma_i'(t)|\,\mathrm dt=0\quad(i=1,2).
\]
因此 \(\gamma'=0\) 在 \(A\) 中几乎处处成立。再对向量映射 \(\gamma\) 使用面积公式的上界，
\[
 \mathcal H^1\bigl(P\cap\gamma(\mathbb R)\bigr)
 =\mathcal H^1\bigl(\gamma(A)\bigr)
 \leq\int_A|\gamma'(t)|\,\mathrm dt=0.
\]
这对任意 Lipschitz 曲线成立，故 \(P\) 纯 \(1\) 维不可整流。
\end{proof}

证明利用的是两个坐标投影都为零测，而非只用“不连通”这种拓扑描述。一般可整流集也可以很不连通，例如一条直线上的正测 Cantor 集。是否有正质量能够落在 Lipschitz 参数片上，才是这里的区分依据。

\subsection{可整流测度}
\label{b1:subsec:190104}

集合之外还可以携带不同的质量密度。为避免把“集中在曲面上”与“沿曲面的面积连续分布”混为一谈，本书明确采用下面的约定。

\begin{definition}\label{b1:def:rectifiable-measure}
若正 Radon 测度 \(\mu\) 满足 \(\mu\ll\mathcal H^k\)，且存在可数 \(k\) 维可整流的 Borel 集 \(E\)，使 \(\mu(\mathbb R^n\setminus E)=0\)，则称 \(\mu\) 为 \(k\) 维\term[kezhengliu-cedu]{可整流测度}[rectifiable measure]。
\end{definition}

绝对连续条件的含义是每个 Borel 的 \(\mathcal H^k\) 零测集都为 \(\mu\) 零测集。这里不要求环境空间上的 \(\mathcal H^k\) 本身是 Radon 测度；当 \(k<n\) 时，它通常不是局部有限的。

\begin{proposition}\label{b1:prop:rectifiable-measure-density}
上述定义等价于存在可数 \(k\) 维可整流 Borel 集 \(E\) 和非负 Borel 函数 \(\theta\)，使
\[
 \mu(B)=\int_{B\cap E}\theta(x)\,\mathrm d\mathcal H^k(x)
 \quad(B\subseteq\mathbb R^n\text{ Borel}).
\]
\(\theta\) 可以取为在 \(E\) 上 \(\mathcal H^k\)-几乎处处有限；若只保留 \(\{\theta>0\}\)，还可令它在承载集上几乎处处为正。简记为 \(\mu=\theta\mathcal H^k|_E\)。
\end{proposition}
\begin{proof}
\(\mathcal H^k|_E\) 为 \(\sigma\) 有限，而 \(\mu\) 是欧氏空间上的 Radon 测度，也为 \(\sigma\) 有限。把第六章的 Radon--Nikodym 定理应用于承载集 \(E\)，即得密度。它可取 Borel 代表；若在正 \(\mathcal H^k\) 测度的集合上为无穷，该集合与某个有界块的交仍有正测度，就会迫使该块具有无穷 \(\mu\) 质量，与局部有限性矛盾。删去 \(\theta=0\) 的部分不损失 \(\mu\) 质量。反过来，所给积分表示直接保证绝对连续性与集中性。
\end{proof}

例如一条线上的 \(\theta\mathcal H^1\) 是带密度的可整流测度，只要每个紧集上的加权质量有限；它的密度不必为整数，也不必等于 \(1\)。单点质量 \(\delta_a\) 虽可集中在任何经过 \(a\) 的直线上，却不是本书意义下的 \(1\) 维可整流测度，因为 \(\mathcal H^1(\{a\})=0\) 而 \(\delta_a(\{a\})=1\)。它应归入零维的计数型几何。

Mattila 的《Rectifiability: A Survey》\cite[定义4.1--4.2及第19页说明]{Mattila2023Rectifiability}对集合采用同样的可数参数覆盖，但对测度有时不附加 \(\mu\ll\mathcal H^k\)。读者对照文献时须保留这一差别；“可整流测度”这个名称本身并不能代替绝对连续假设。本书保留该假设，是为了使面积密度、切向积分和后文的测度密度结论使用同一约定。

% !TeX root = ../../Book1.tex
\section{近似切空间}
\label{b1:sec:1902}

Lipschitz 参数化在几乎每一点有线性微分，但一个可整流集可能由许多相交的参数像组成。要把参数空间中的微分解释成集合自身的切平面，既要看清单个片的局部形状，也要控制其余各片进入同一个小球的质量。本节先用正上密度与锥外质量定义近似切平面，再把单位密度的平面缩放极限作为更强结论单列。这样既与通常的 GMT 定义对齐，也不丢失 Simon 的《Lectures on Geometric Measure Theory》\cite[11.2 DEFINITION, 11.3 REMARK, 11.4 DEFINITION, 11.5 REMARK, and 11.6 THEOREM]{Simon1983Lectures}所强调的缩放测度信息。本节直接从 Lipschitz 片和已证的面积公式出发，不另引入光滑逼近定理。

% 来源：实读 Simon 1983 印刷第59--63页，第11.2--11.7节；单位密度缩放结论及片间转移思路据此，
% 标准锥定义与重数区分按本轮定义审计单列，
% 其 C^1 分片在本节改为第17章双 Lipschitz 分片、加权面积公式与 Lebesgue 点的完整论证。
% 另实读 Fremlin 作者版 chap47.tex 第474O--P：该处 Federer exterior normal
% 针对实体集合的半空间体积逼近，不是本节一般余维的近似法空间，未据此移用结论。

\subsection{近似切平面}
\label{b1:subsec:190201}

以下先固定 \(1\leq k\leq n\)，并设 \(E\subseteq\mathbb R^n\) 为 Borel 集，限制测度 \(\mu=\mathcal H^k|_E\) 局部有限。这里 \(\mathcal H^k|_E(A)=\mathcal H^k(E\cap A)\)，不是把整个空间上的 \(\mathcal H^k\) 假定为局部有限。对 \(k<n\)，后一个假定通常不成立。

\begin{definition}\label{b1:def:approximate-tangent-plane}
设 \(x\in E\)，\(T\subseteq\mathbb R^n\) 为 \(k\) 维线性子空间，并且
\[
 \limsup_{r\downarrow0}
 \frac{\mathcal H^k\bigl(E\cap\overline B(x,r)\bigr)}
      {\omega_kr^k}>0.
\]
若对每个 \(\varepsilon>0\) 都有
\begin{equation}\label{b1:eq:approximate-tangent-cone}
 \lim_{r\downarrow0}\frac1{r^k}\mathcal H^k
 \left(\left\{\,y\in E\cap\overline B(x,r):
  \operatorname{dist}(y-x,T)\geq\varepsilon|y-x|\,\right\}\right)=0,
\end{equation}
则称 \(T\) 为 \(E\) 在 \(x\) 处的\term[jinsi-qiepingmian]{近似切平面}[approximate tangent plane]。在以下几乎处处唯一的情形中，记为 \(\operatorname{Tan}^k(E,x)\)；对应的仿射平面为 \(x+\operatorname{Tan}^k(E,x)\)。
\symidx{\(\operatorname{Tan}^k(E,x)\)：可整流集的近似切平面}
\end{definition}

这个定义只要求小球中的主阶质量逐渐集中到 \(T\) 的任意细锥内，并不预先指定密度为 \(1\)。改变 \(E\) 的一个 \(\mathcal H^k\) 零测子集不会改变密度与锥条件，但是否把某个点列为 \(E\) 的元素，仍取决于所选集合。

\cref{b1:fig:approximate-tangent-cone}示意锥条件的几何含义。随着观察半径缩小，集合允许有少量点落在锥外，但这部分的归一化面积必须趋于零。
\begin{figure}[htbp]
  \centering
  \begin{tikzpicture}[scale=1.15]
    \fill[gray!13] (0,0)--(2,0.9)--(2,-0.9)--cycle;
    \fill[gray!13] (0,0)--(-2,0.9)--(-2,-0.9)--cycle;
    \draw[dashed] (-2,-0.9)--(0,0)--(2,0.9);
    \draw[dashed] (-2,0.9)--(0,0)--(2,-0.9);
    \draw[thick] (-2.15,0)--(2.15,0) node[right] {$x+T$};
    \draw[thick,domain=-1.85:1.85,samples=70] plot (\x,{0.22*\x*\x});
    \draw (0,0) circle (2);
    \fill (0,0) circle (1.5pt) node[below left] {$x$};
    \node at (1.18,-0.43) {细锥};
    \node[above] at (0,1.98) {$B(x,r)$};
  \end{tikzpicture}
  \caption[近似切平面]{近似切平面。灰色区域表示围绕仿射平面 \(x+T\) 的细锥；锥外质量在缩放后可以忽略。}
  \label{b1:fig:approximate-tangent-cone}
\end{figure}

后文实际证明的结论更强：对每个 \(\varphi\in C_c(\mathbb R^n)\) 都有
\begin{equation}\label{b1:eq:approximate-tangent-blowup}
 \lim_{r\downarrow0}\frac1{r^k}
 \int_E\varphi\left(\frac{y-x}{r}\right)\,\mathrm d\mathcal H^k(y)
 =\int_T\varphi(z)\,\mathrm d\mathcal H^k(z).
\end{equation}
这称为在 \(x\) 处具有方向 \(T\) 的单位密度平面缩放极限。它是一项带有单位重数的加强性质，不是“近似切平面”的第二个同义定义；一般带权测度的缩放极限将在\secref{b1:sec:1904}另行讨论。局部有限性保证上式左端的每次积分有限。

\begin{proposition}\label{b1:prop:approximate-tangent-properties}
若 \(E\) 在 \(x\) 处满足\eqref{b1:eq:approximate-tangent-blowup}，则满足该缩放条件的平面 \(T\) 唯一，并且
\begin{equation}\label{b1:eq:rectifiable-unit-density}
 \lim_{r\downarrow0}
 \frac{\mathcal H^k\bigl(E\cap\overline B(x,r)\bigr)}
      {\omega_k r^k}=1.
\end{equation}
开球也满足同一极限，而且 \(T\) 是\cref{b1:def:approximate-tangent-plane}意义下的近似切平面。
\end{proposition}
\begin{proof}
若 \(T_1,T_2\) 为两个不同的候选平面，取 \(z\in T_1\setminus T_2\)，再取一个与 \(T_2\) 不交的小球 \(B(z,\rho)\)。可以选取非负连续函数 \(\varphi\)，支集包含于这个球，并在 \(z\) 的某个邻域严格为正。它在 \(T_1\) 上的积分为正，在 \(T_2\) 上的积分为零，与同一个缩放积分具有这两个极限矛盾。

为求球的质量，固定 \(0<\delta<1\)，取径向连续函数 \(\varphi_-,\varphi_+\)，满足
\[
 {\bf1}_{\overline B(0,1-\delta)}
 \leq\varphi_-\leq{\bf1}_{B(0,1)}
 \leq{\bf1}_{\overline B(0,1)}
 \leq\varphi_+\leq{\bf1}_{B(0,1+\delta)}.
\]
由\eqref{b1:eq:approximate-tangent-blowup}，开球与闭球的归一化质量下极限至少为 \((1-\delta)^k\)，上极限至多为 \((1+\delta)^k\)。这里使用了平面上 \(\mathcal H^k\) 与 \(k\) 维 Lebesgue 测度的一致性。令 \(\delta\downarrow0\) 即得密度结论。

最后固定 \(\varepsilon,\delta>0\)。紧集
\[
 F_{\varepsilon,\delta}
 =\{\,z:\delta\leq|z|\leq1,\,
              \operatorname{dist}(z,T)\geq\varepsilon|z|\,\}
\]
与 \(T\) 不交，故可取非负的紧支集连续函数 \(\chi\)，在这个紧集上等于 \(1\)，并在 \(T\) 的一个邻域等于零。缩放积分条件说明落在该环形区域内的锥外质量除以 \(r^k\) 趋于零。剩余部分在 \(\overline B(x,\delta r)\) 中，其质量除以 \(r^k\) 的极限为 \(\omega_k\delta^k\)。再令 \(\delta\downarrow0\)，便得到\eqref{b1:eq:approximate-tangent-cone}。单点的 \(\mathcal H^k\) 测度为零，因而 \(y=x\) 不影响结论。
\end{proof}

锥外质量的消失说明，少量偏离切平面的点可以保留，但它们在当前尺度下不能占有正比例的面积。这个条件不要求小球内的整个集合都靠近平面，也不要求集合在一点的邻域内恰是一张图。它也不强迫单位密度。例如令
\[
 E=\{\,(t,t^2):|t|<1\,\}\cup
   \{\,(t,-t^2):|t|<1\,\}.
\]
两条曲线在原点相切。对任意锥角，充分靠近原点的两条曲线都落在水平轴的锥内，所以水平轴是标准近似切线；但两支分别贡献一份长度，故原点的密度为 \(2\)，缩放测度趋于水平轴上 \(2\mathcal H^1\)，而不是\eqref{b1:eq:approximate-tangent-blowup}中的单位重数极限。这正是两层陈述不能逐点混同的原因。

\subsection{近似法空间}
\label{b1:subsec:190202}

\begin{definition}\label{b1:def:approximate-normal-space}
若 \(\operatorname{Tan}^k(E,x)\) 存在，则定义 \(E\) 在 \(x\) 处的\term[jinsi-fakongjian]{近似法空间}[approximate normal space]为
\[
 \operatorname{Nor}^k(E,x)
 \coloneqq\operatorname{Tan}^k(E,x)^\perp.
\]
其中的向量称为\term[jinsi-faxiangliang]{近似法向量}[approximate normal vector]；长度为 \(1\) 的向量称为单位近似法向量。
\symidx{\(\operatorname{Nor}^k(E,x)\)：可整流集的近似法空间}
\end{definition}

在余维 \(1\) 的情形，单位法向量恰有两个互为相反数的选择。仅从集合本身的无向面积测度出发，没有指定其中哪一个为“外法向量”。余维大于 \(1\) 时，法空间的维数为 \(n-k\)，更不能用唯一的单位向量代替它；当 \(k=n\) 时，法空间只有零向量。

例如，设 \(g:\mathbb R^k\rightarrow\mathbb R\) 为 Lipschitz 函数，\(\psi(a)=(a,g(a))\)。在 \(g\) 可微且图的近似切平面等于 \(\operatorname{im}D\psi(a)\) 的点，切向量与一个单位法向量分别为
\[
 (v,Dg(a)v),\quad
 \nu(a)=\frac{(-\nabla g(a),1)}
                    {\sqrt{1+|\nabla g(a)|^2}}.
\]
两者内积为零，这直接验证了法向量公式。切平面与参数微分几乎处处相容的事实将在本节后半部证明。此处的法向量针对图上的面积；一个实体集合由哪一侧占据半空间，是另一个需要补充内外信息的问题。

\subsection{密度与切空间}
\label{b1:subsec:190203}

先把问题限于一个没有自交的参数片。面积公式负责把像上的积分拉回参数空间，Lebesgue 点则保证缩小邻域中的 Jacobi 因子平均趋于它在中心的值。

\begin{lemma}\label{b1:lem:bilipschitz-piece-blowup}
设 \(F:\mathbb R^k\rightarrow\mathbb R^n\) 为 Lipschitz 映射，\(D\subseteq\mathbb R^k\) 为 Borel 集，且存在 \(c>0\) 使
\[
 |F(u)-F(v)|\geq c|u-v|\quad(u,v\in D).
\]
记 \(M=F(D)\)，并设 \(DF\) 在 \(D\) 的几乎每一点存在且秩为 \(k\)。则对 \(\mathcal H^k\)-几乎处处的 \(x=F(a)\in M\)，集合 \(M\) 的缩放积分满足\eqref{b1:eq:approximate-tangent-blowup}，其中 \(T=\operatorname{im}DF(a)\)。
\end{lemma}
\begin{proof}
采用\secref{b1:sec:1702}的 Borel Jacobi 因子，置 \(h={\bf1}_D J_kF\)。它是有界可测函数。删去一个参数零测集后，可同时要求 \(F\) 在 \(a\in D\) 处可微、\(A=DF(a)\) 满秩，以及 \(a\) 为 \(h\) 的 Lebesgue 点，且 \(h(a)=J_k(A)\)。被删集合的像也是 \(\mathcal H^k\) 零测集。

像 \(M\) 由\cref{b1:lem:lipschitz-image-hausdorff-measurable}可测。其局部质量有限：若 \(M\cap\overline B(x,R)\) 非空，选取其中一点的原像，则所有其他原像均落在一个半径 \(2R/c\) 的参数球内；Lipschitz 测度估计给出有限的 \(\mathcal H^k\) 质量。因此本引理的紧支集积分有意义。

固定上述好点 \(a\)，令 \(x=F(a)\)。由单射情形的加权面积公式与参数空间的平移缩放，
\[
 \begin{aligned}
 &\frac1{r^k}\int_M
       \varphi\left(\frac{y-x}{r}\right)\,\mathrm d\mathcal H^k(y)\\
 &\hspace{1em}=
 \int_{\mathbb R^k}
   \varphi\left(\frac{F(a+rv)-F(a)}r\right)h(a+rv)\,\mathrm dv.
 \end{aligned}
\]
若 \(\operatorname{supp}\varphi\subseteq\overline B(0,R)\)，被积函数非零时必有 \(a+rv\in D\)，进而 \(|v|\leq R/c\)。满秩矩阵 \(A\) 有最小伸缩 \(\sigma>0\)，所以 \(\varphi(Av)\) 非零时也有 \(|v|\leq R/\sigma\)。选取大于这两个界的常数 \(L\)，以后只在 \(B(0,L)\) 上积分。

可微性给出
\[
 \sup_{|v|\leq L}
 \left|\frac{F(a+rv)-F(a)}r-Av\right|\longrightarrow0.
\]
而 Lebesgue 点性质给出
\[
 \int_{B(0,L)}|h(a+rv)-J_k(A)|\,\mathrm dv
 =\frac1{r^k}\int_{B(a,Lr)}|h(u)-h(a)|\,\mathrm du
 \longrightarrow0.
\]
利用 \(\varphi\) 的一致连续性和有界性，前一个积分因而趋于
\[
 J_k(A)\int_{\mathbb R^k}\varphi(Av)\,\mathrm dv
 =\int_{\operatorname{im}A}\varphi\,\mathrm d\mathcal H^k,
\]
末步是线性面积公式。所选好点只取决于可微性与 \(h\) 的 Lebesgue 点性质，不取决于测试函数；故在这些点对全部 \(\varphi\in C_c\) 同时成立。
\end{proof}

从单片走向整个集合时，不能仅说“各片不交，所以其他片没有贡献”：互不相交的片仍可在任意小球中彼此逼近。下一证明用奇异测度微分定理控制这种贡献。

\begin{theorem}\label{b1:thm:rectifiable-density-tangent}
设 \(E\subseteq\mathbb R^n\) 为 Borel 的可数 \(k\)-可整流集，且 \(\mathcal H^k|_E\) 局部有限。则对 \(\mathcal H^k\)-几乎每个 \(x\in E\)，存在唯一的 \(k\) 维平面 \(T_x\) 使单位密度平面缩放极限\eqref{b1:eq:approximate-tangent-blowup}成立。特别地，\(T_x\) 是 \(E\) 在 \(x\) 处的近似切平面，记为 \(\operatorname{Tan}^k(E,x)\)，并且
\[
 \lim_{r\downarrow0}
 \frac{\mathcal H^k\bigl(E\cap\overline B(x,r)\bigr)}
      {\omega_k r^k}=1
 \quad\text{对 }\mathcal H^k\text{-几乎处处的 }x\in E.
\]
\end{theorem}
\begin{proof}
按可数可整流性的定义，用可数个全空间 Lipschitz 映射 \(F_i:\mathbb R^k\rightarrow\mathbb R^n\) 的像覆盖 \(E\)，只遗漏一个 \(\mathcal H^k\) 零测集。若原始参数域不是全空间，先用\secref{b1:sec:1601}的逐坐标延拓。

对每个 \(F_i\)，不可微集的像为 \(\mathcal H^k\) 零测集，秩亏集的像由\cref{b1:lem:rank-deficient-image-null}也为零测集。由\cref{b1:lem:linear-approximation-pieces}，其满秩集可分为可数个 Borel 片，使每片上 \(F_i\) 都是双 Lipschitz 映射。具体地，把对应误差取为小于单射矩阵最小伸缩的一半，就得到严格的距离下界。

将每个参数片与 \(F_i^{-1}(E)\) 相交，再用 Lebesgue 内正则性在其有界部分中选取可数紧集，使遗漏参数集为零测集。全部这些紧集的像记为 \(Q_1,Q_2,\ldots\)；它们是包含于 \(E\) 的紧集，并在忽略零测集后覆盖 \(E\)。依次相减，得到两两不交的 Borel 片
\[
 M_j=Q_j\setminus\bigcup_{\ell<j}Q_\ell.
\]
若 \(Q_j=F_{i_j}(K_j)\)，则
\[
 D_j=K_j\cap F_{i_j}^{-1}(M_j)
\]
为 Borel 集，\(F_{i_j}|_{D_j}\) 双 Lipschitz，且其微分在 \(D_j\) 上满秩。由上一引理，\(M_j\) 在自身几乎每一点的缩放极限为一个 \(k\) 维平面，并在这些点有密度 \(1\)。

现在置
\[
 \mu=\mathcal H^k|_E,\quad
 \nu_j=\mathcal H^k|_{M_j},\quad
 \sigma_j=\mathcal H^k|_{E\setminus M_j}.
\]
这些都是局部有限的 Borel 测度，因而是本书约定的 Radon 测度。为明确正则性的来源，把它们限制到 \(\overline B(0,N)\)，应用\cref{b1:lem:polish-probability-radon}中有限 Borel 测度的紧内正则性；再令 \(N\) 递增，便得到所有 Borel 集的紧内正则性。

由于 \(\sigma_j\perp\nu_j\)，由\cref{b1:thm:radon-singular-differentiation}，在 \(\nu_j\)-几乎每一点有
\[
 \frac{\sigma_j\bigl(\overline B(x,r)\bigr)}
      {\nu_j\bigl(\overline B(x,r)\bigr)}
 \longrightarrow0.
\]
在 \(M_j\) 的密度为 \(1\) 的点，分母等价于 \(\omega_k r^k\)，故
\[
 \sigma_j\bigl(\overline B(x,r)\bigr)=o(r^k).
\]
这一论证使用中心闭球，并未假设任何球边界对这些测度为零测集。

对支集在 \(\overline B(0,R)\) 内的测试函数，整个 \(E\) 与单片 \(M_j\) 的缩放积分之差的绝对值不超过
\[
 \|\varphi\|_\infty\,
 \frac{\sigma_j\bigl(\overline B(x,Rr)\bigr)}{r^k},
\]
它趋于零。因此 \(M_j\) 的平面缩放极限也是 \(E\) 的缩放极限。对所有 \(j\) 取可数并，便在 \(E\) 的几乎每一点建立了单位密度平面缩放极限；其唯一性、近似切平面性质与密度由\cref{b1:prop:approximate-tangent-properties}给出。
\end{proof}

局部有限性在这里控制的是所有片合起来的质量，保证片间余量可以作为 Radon 测度来微分。单个参数像的可数覆盖本身不能替代这个假设。

即使密度在某一点恰好为 \(1\)，该点仍可能没有近似切平面。取平面内两条互相垂直的正半轴
\[
 E=\{\,te_1:t\geq0\,\}\cup\{\,te_2:t\geq0\,\}.
\]
在原点有 \(\mathcal H^1\bigl(E\cap B(0,r)\bigr)=2r=\omega_1r\)。但任意一条过原点的直线都至少与一条半轴成正角；选择足够小的锥角后，该半轴在球中的长度始终为 \(r\)，违反\eqref{b1:eq:approximate-tangent-cone}。原点因而没有近似切线。密度记录质量总数，不能独自辨别它分布在哪些方向。

\subsection{Lipschitz 图表示}
\label{b1:subsec:190204}

一般的 Lipschitz 参数像在全局可以折叠或相交。要在集合上定义微分，采用没有自交且坐标投影可逆的图更方便。\secref{b1:sec:1703}的近线性分片已经提供了这种图，只需把片上定义的函数延拓到整个平面。

\begin{proposition}\label{b1:prop:rectifiable-lipschitz-graphs}
在\cref{b1:thm:rectifiable-density-tangent}的假设下，存在可数个 \(k\) 维线性平面 \(P_j\)、Lipschitz 映射 \(g_j:P_j\rightarrow P_j^\perp\)，以及两两不交的 Borel 集
\[
 M_j\subseteq E\cap\Gamma_j,\quad
 \Gamma_j=\{\,z+g_j(z):z\in P_j\,\},
\]
使 \(\mathcal H^k\bigl(E\setminus\bigcup_jM_j\bigr)=0\)。记
\[
 \psi_j(z)=z+g_j(z),\quad A_j=\psi_j^{-1}(M_j).
\]
则 \(A_j\) 为 Borel 集，\(\psi_j|_{A_j}\) 双 Lipschitz；对参数平面上几乎处处的 \(a\in A_j\)，有
\begin{equation}\label{b1:eq:graph-approximate-tangent}
 \operatorname{Tan}^k\bigl(E,\psi_j(a)\bigr)
 =\operatorname{im}D\psi_j(a).
\end{equation}
这里参数平面的零测集与 \(M_j\) 中的 \(\mathcal H^k\) 零测集通过 \(\psi_j\) 互相对应。
\end{proposition}
\begin{proof}
先设 \(k<n\)。在一个满秩参数片 \(D\) 上，\secref{b1:sec:1703}的分片估计给出单射矩阵 \(T\) 与 \(\eta>0\)，使
\[
 |F(v)-F(u)-T(v-u)|\leq\eta|v-u|
 \quad(u,v\in D),
\]
并可要求 \(\eta<\sigma/2\)，其中 \(\sigma\) 是 \(T\) 的最小伸缩。令 \(P=\operatorname{im}T\)，以 \(p,q\) 分别表示到 \(P,P^\perp\) 的正交投影。于是
\[
 \begin{aligned}
 |pF(v)-pF(u)|&\geq(\sigma-\eta)|v-u|,\\
 |qF(v)-qF(u)|&\leq\eta|v-u|.
 \end{aligned}
\]
第一式使 \(p\) 在 \(F(D)\) 上单射，因此
\[
 h\bigl(pF(u)\bigr)=qF(u)
\]
在 \(pF(D)\) 上良定义。两式相除表明 \(h\) 的 Lipschitz 常数不超过 \(\eta/(\sigma-\eta)\)。在 \(P^\perp\) 选定标准正交基，逐坐标应用\cref{b1:thm:mcshane-extension}，便把 \(h\) 延拓为全平面 \(P\) 上的 Lipschitz 映射 \(g\)。延拓的向量 Lipschitz 常数至多额外乘以 \(\sqrt{n-k}\)，其有限性已经足够。所得完整图包含 \(F(D)\)。

把这一构造用于可数参数像的全部满秩片；忽略已经证明像为零测集的部分，得到可数图 \(\Gamma_j\) 覆盖 \(E\) 的几乎全部。每张完整图是闭集：若 \(z_\ell+g_j(z_\ell)\) 收敛，投影得到 \(z_\ell\) 收敛，连续性再给出极限仍在图上。故可置
\[
 M_j=(E\cap\Gamma_j)\setminus\bigcup_{\ell<j}\Gamma_\ell,
\]
它们 Borel、两两不交，并满足所需覆盖。连续性保证 \(A_j\) Borel。正交投影是 \(\psi_j\) 的逆，且
\[
 |a-b|\leq|\psi_j(a)-\psi_j(b)|
 \leq\sqrt{1+\operatorname{Lip}(g_j)^2}\,|a-b|,
\]
所以参数化双 Lipschitz。

在 \(g_j\) 可微的点，\(D\psi_j(v)=v+Dg_j(v)\) 的 \(P_j\) 分量就是 \(v\)，故它满秩。对 \(A_j\) 应用\cref{b1:lem:bilipschitz-piece-blowup}，再像前一定理那样对 \(\mathcal H^k|_{E\setminus M_j}\) 与 \(\mathcal H^k|_{M_j}\) 使用奇异测度微分，得到\eqref{b1:eq:graph-approximate-tangent}。参数零测集由 Lipschitz 性映为面积零测集；反向对零测子集使用逆投影的 Lipschitz 性即可。

当 \(k=n\) 时只需取 \(P_1=\mathbb R^n\)、\(g_1=0\)、\(M_1=E\)。图参数化为恒等映射，结论由同一个单片论证成立。
\end{proof}

这些图给出的只是可数覆盖和片上的坐标，并不说明 \(E\) 在几乎每一点的整个邻域内等于一张图。邻近的其他片仍可存在；证明已经说明，它们在典型点的面积贡献是低阶量。

为后续积分，还可将切平面的可测性写成矩阵形式。对上述图的好点，令 \(A=D\psi_j(a)\)，则到近似切平面的正交投影矩阵为
\[
 \Pi\bigl(\psi_j(a)\bigr)=A(A^{\mathsf T}A)^{-1}A^{\mathsf T}.
\]
Borel 微分代表、连续的逆投影坐标及矩阵求逆说明，这个矩阵在各 Borel 片上有 Borel 代表。把各片的例外集包在一个 Borel 零测集中，并在那里指定一个固定 \(k\) 维平面的投影，就得到整个 \(E\) 上的 Borel 代表，且几乎处处确为到 \(\operatorname{Tan}^k(E,x)\) 的投影。下一节因此可以在切平面上定义微分与 Jacobi 因子，而不必要求某个环境延拓恰好在 \(\mathcal H^k\)-几乎每一点可微。

若另考虑 \(k=0\)，局部有限的 \(\mathcal H^0|_E\) 意味着每个有界区域只含有限个点。每个 \(x\in E\) 因而在 \(E\) 中孤立，其缩放测度趋于原点的单位点质量；对应的切空间为 \(\{0\}\)，法空间为 \(\mathbb R^n\)。以上需要 \(r^k\) 质量估计的论证主要服务于正维数情形。

% !TeX root = ../../Book1.tex
\section{可整流集上的面积公式}
\label{b1:sec:1903}

% TODO: 撰写本节正文。

\subsection{参数化}
\label{b1:subsec:190301}

待填充。

\subsection{变量代换}
\label{b1:subsec:190302}

待填充。

\subsection{切向 Jacobi 行列式}
\label{b1:subsec:190303}

待填充。

\subsection{可整流测度}
\label{b1:subsec:190304}

待填充。

% !TeX root = ../../Book1.tex
\OptionalSection{密度与可整流性}{b1:sec:1904}

本章主线已经从参数化推出切平面、集合密度与带权可整流测度的密度。现在反过来问：如果不知道参数化，只观察小球内的质量，能否判断集合或测度可整流？这一逆问题远比顺向计算困难。本选读节先证明可整流／纯不可整流分解，再陈述作为深层输入的 Preiss 定理，最后把集合的密度判据作为推论。深层几何步骤明确交由文献，不能把它们看作前面面积公式证明的一部分。

% 实读：Mattila2023Rectifiability 出版社OA正式全书§4.3--4.4，尤其pp.23--25定理4.10--4.11及证明思路。
% Preiss1987Geometry 的Annals官方页核验作者、卷期页码DOI，但未获得其107页原文的完整证明。
% 密度判据的逆向以Preiss为深层输入；Preiss只给明确标出的proofsketch。
% 带权密度已移入19.3主线；rectifiable/pure分解为本书从有限质量上确界完整推导。

\subsection{Federer 型结构定理}
\label{b1:subsec:190403}

无论密度极限是否存在，有限 \(\mathcal H^k\) 测度的 Borel 集总能分出可被参数片捕捉的部分与完全剩余的部分。这一分解不需要 Preiss 定理。

\begin{theorem}[可整流与纯不可整流分解]\label{b1:thm:rectifiable-pure-decomposition}
设 \(E\subseteq\mathbb R^n\) 是 Borel 可测的，且 \(\mathcal H^k(E)<\infty\)。则可以写成不交并
\[
 E=R\mathbin{\dot\cup}P,
\]
其中 \(R\) 是 Borel 的且可数 \(k\) 维可整流，\(P\) 是 Borel 的且纯 \(k\) 维不可整流。这样的分解在 \(\mathcal H^k\) 零测集意义下唯一。
\end{theorem}
\begin{proof}
在所有可数 \(k\) 维可整流的 Borel 子集 \(F\subseteq E\) 中，取其测度的上确界
\[
 a=\sup\{\mathcal H^k(F):F\subseteq E\text{ 为这样的子集}\}.
\]
该数有限。取 \(F_j\) 使其测度趋于 \(a\)，并令 \(R=\bigcup_jF_j\)。可数并仍可整流，所以 \(\mathcal H^k(R)\leq a\)；又因它包含每个 \(F_j\)，有 \(\mathcal H^k(R)=a\)。置 \(P=E\setminus R\)。

任取 Lipschitz 映射 \(f:\mathbb R^k\to\mathbb R^n\)。其全空间像为 Borel 集，故 \(A=P\cap f(\mathbb R^k)\) 为 Borel 的可整流集。如果 \(\mathcal H^k(A)>0\)，则不交并 \(R\cup A\) 仍是 \(E\) 的可整流 Borel 子集，但
\[
 \mathcal H^k(R\cup A)=a+\mathcal H^k(A)>a,
\]
与上确界的定义矛盾。因此 \(P\) 纯不可整流。

若还有 \(E=R'\mathbin{\dot\cup}P'\) 满足同样条件，则 \(R\setminus R'\subseteq R\cap P'\) 为零测集，因为 \(R\) 可整流而 \(P'\) 纯不可整流。交换两组分解，得到 \(R'\setminus R\) 也零测。故 \(R\) 与 \(R'\) 的对称差零测，补集 \(P,P'\) 也只差零测集。
\end{proof}

若 \(E\) 只有 \(\sigma\) 有限测度，可先分成可数个两两不交的有限测度 Borel 块，逐块分解后分别取并。可整流部分的可数并仍可整流，纯不可整流部分与每个参数像的交则是可数个零测集之并，所以仍纯不可整流。唯一性采用同一交集论证，不需要对无穷质量作相减。

这一结构体现了 Besicovitch 的平面理论及 Federer 的高维推广中最基本的区分：参数化所能识别的部分可以集中保留，剩下的部分与每张参数片都只零测相交。完整的 Federer 型结构理论还研究投影、近似切平面与纯不可整流部分之间的关系。那些投影判据需要另外指定“几乎每个方向”所用的测度，不应从上面的上确界分解自行推出。可进一步参阅 Federer 的《Geometric Measure Theory》\cite{Federer1996Geometric}及 Mattila 的《Rectifiability: A Survey》\cite[Section 4]{Mattila2023Rectifiability}。

\subsection{Preiss 定理概览}
\label{b1:subsec:190404}

现在陈述从密度重建可整流结构所需的深层结果，即\term[preiss-dingli]{Preiss 定理}[Preiss's theorem]。与集合测度不同，一般 Radon 测度的密度不必等于 \(1\)，因此正确假设是密度存在、正且有限。

\begin{theorem}[Preiss]\label{b1:thm:preiss-rectifiability}
设 \(\mu\) 为 \(\mathbb R^n\) 上的正 Radon 测度，\(1\leq k\leq n\) 为整数。如果
\[
 0<\lim_{r\downarrow0}\frac{\mu(B(x,r))}{\omega_kr^k}<\infty
 \quad\text{对 }\mu\text{-几乎每个 }x,
\]
则 \(\mu\) 是本书意义下的 \(k\) 维可整流测度。反过来，每个这样的可整流测度都满足上述密度条件。
\end{theorem}

定理的顺向几何结论来源于 Preiss 的《Geometry of Measures in \(\mathbf R^n\): Distribution, Rectifiability, and Densities》\cite{Preiss1987Geometry}。本书所述版本可与 Mattila 的《Rectifiability: A Survey》\cite[Theorem 4.11 and the following discussion]{Mattila2023Rectifiability}核对；其中从有限上密度得到 \(\mu\ll\mathcal H^k\)，所以也满足本书较强的测度约定。反向已经由\cref{b1:prop:rectifiable-weight-density}完整证明。下面只解释正向证明的组织，不尝试用几段话替代原论文的长证明。

\begin{proofsketch}
先看缩放。对 \(x\in\mathbb R^n\)、\(r>0\)，令 \(T_{x,r}(z)=(z-x)/r\)。把 \(\mu\) 推前并重新调整质量后，考察
\[
 c_j(T_{x,r_j})_\#\mu,
 \quad c_j>0,\quad r_j\downarrow0,
\]
在 \(C_c(\mathbb R^n)\) 测试下的非零 Radon 极限。这样的极限称为 \(\mu\) 在 \(x\) 处的\term[qiecedu]{切测度}[tangent measure]。它是一个测度，而不是一张线性平面；其收敛是\chapref{b1:ch:12}的淡收敛。

正的有限密度给缩放测度提供局部质量控制。以 \(c_j=r_j^{-k}\) 为例，每个固定半径球的质量保持有界，原点附近又不会全部消失。这使提取非零局部极限成为可能。仅有这一紧性还不够，重要的是把原测度在几乎处处的密度信息转移到这些极限的全部支撑点，得到统一的球质量关系。

接下来需要研究满足该关系的测度的结构。这里有一个不可省略的困难：当维数较高时，具有同一幂次球质量的测度不一定集中在一张平面上。因此不能在得到这种极限后直接宣布“切测度是平面的”。Preiss 的证明还使用切测度族的结构和进一步缩放，从中取得足够的平坦性，再回到原测度，构造可数个捕捉其质量的参数片。

本概要省略了局部极限的系统构造、统一球质量的传递、非平坦极限的排除机制以及最后的可整流性恢复；这些正是定理的深层部分，而非例行估计。Mattila 的《Rectifiability: A Survey》\cite[Sections 4.3 Tangent Measures and 4.4 Densities]{Mattila2023Rectifiability}给出这些步骤之间的关系和进一步文献，完整论证须读 Preiss 原文。这里没有把任意一个切测度的存在误当成已经得到可整流性。
\end{proofsketch}

切平面与切测度的区别也可以在最熟悉的例子上看出。若 \(\mu=\theta\mathcal H^k|_V\)，其中 \(V\) 为平面、\(\theta>0\) 为常数，则在 \(x\in V\) 处，\(r^{-k}(T_{x,r})_\#\mu=\theta\mathcal H^k|_{V-x}\)。平面给出方向，测度另外保留质量常数；更一般的缩放极限还可能保留重数和非平坦形状。

Preiss 定理的结论是测度意义下可数参数化，并不把承载集变成处处光滑的流形，不保证全局一张图，也不强迫密度为整数。它说明，小球质量在全部小尺度上稳定到一个正有限幂律时，欧氏几何中会出现可整流结构。

\subsection{可整流性的密度判据}
\label{b1:subsec:190402}

顺向结论只用了微分、面积与测度微分。逆向则需要从小球质量中重建方向；有了上一小节的 Preiss 定理，现在可以把这一步作为正式推论写出。

\begin{theorem}[可整流性的密度判据]\label{b1:thm:rectifiability-density-criterion}
设 \(E\subseteq\mathbb R^n\) 是 Borel 可测的，\(1\leq k\leq n\)，且 \(\mathcal H^k(E)<\infty\)。下列条件等价：
\begin{equivalences}
\item\label{b1:item:density-criterion-rectifiable}\(E\) 可数 \(k\) 维可整流；
\item\label{b1:item:density-criterion-one}\(\Theta^k(\mathcal H^k|_E,x)=1\) 对 \(\mathcal H^k\)-几乎每个 \(x\in E\) 成立；
\item\label{b1:item:density-criterion-exists}\(\Theta^k(\mathcal H^k|_E,x)\) 对 \(\mathcal H^k\)-几乎每个 \(x\in E\) 存在。
\end{equivalences}
\end{theorem}
\begin{proof}[由 Preiss 定理推出]
第一项推出第二项由\cref{b1:thm:rectifiable-density-tangent}给出，第二项当然推出第三项。若第三项成立，由\chapref{b1:ch:15}的 Hausdorff 上密度估计，
\[
 2^{-k}\leq\Theta^{*k}(\mathcal H^k|_E,x)\leq1
\]
在 \(E\) 中几乎处处成立。因此一旦密度极限存在，它必正且有限。将\cref{b1:thm:preiss-rectifiability}应用于有限 Radon 测度 \(\mu=\mathcal H^k|_E\)，得到 \(\mu\) 集中在可数个 Lipschitz 像上。因而这些像在 \(E\) 中遗漏的 \(\mathcal H^k\) 测度为零，这正是第一项。
\end{proof}

这里已经完整写出各条件之间的推导，但逆向所调用的 Preiss 定理并未在本书内证明。Mattila 的《Rectifiability: A Survey》\cite[Theorems 4.10--4.11]{Mattila2023Rectifiability}将这两层结果相邻表述，便于核对有限测度、几乎处处及欧氏空间等假设。前 3 节关于参数化和切向面积的结论都不反向依赖这条深层判据。

单有正的有限上密度不够：\chapref{b1:ch:15}已经证明这一性质对任意有限 \(\mathcal H^k\) 测度的 Borel 集几乎处处成立，而四角 Cantor 集提供了纯不可整流的例子。因此新的要求不是某个尺度上存在面积，而是所有趋零尺度上的归一化质量具有同一个极限。它把无穷多个尺度联系起来，不能用某一列半径上的收敛替代。至此，本编从覆盖与微分走到了低维集合的局部形状；下一编将用弱导数，并在 Sobolev 空间和有界变差函数空间中研究函数，届时这些几何工具会重新出现在迹、边界和有限周长问题中。

% 本章小结（不计入编号）
\begin{chaptersummary}
可数可整流性允许可数参数片和面积零测的余项，纯不可整流性则排除了与任何参数片的正测相交。四角 Cantor 集说明，一个集合可以有正且有限的长度，却没有可由曲线捕捉的正长度部分。有限面积集合总能在零测意义下唯一分成这两类部分。

局部有限的可整流集合在几乎处处具有单位重数的平面缩放极限。证明先把一张双 Lipschitz 片上的积分拉回参数域，再以测度微分控制其余片的质量。图参数化随后给出切向微分，重叠参数域的密度点保证它与坐标和延拓的选择无关。切向面积公式计算的是面积伸缩后的质量；原质量的推前还须保留被临界映射压缩的部分。

密度的顺向结论已在本章证明；由正有限密度反推可整流性则属于 Preiss 定理的深层内容，这里只概述其方法与边界。几何测度论的这一组工具为后面的弱导数、Sobolev 和 BV 理论准备了低维集合及积分的语言。
\end{chaptersummary}


\BookExercises

% 第十九章习题来源与设计说明。
% 八题均为依据本章及前文工具组织的自拟变式；不冒认教材题号。
% 不以四角 Cantor 集的纯不可整流性或尚未证明的 Preiss 定理作为解题前置。

\begin{exercise}\label{b1:ex:19-dense-parallel-planes}
设 \(k\ge1\)，在 \(\mathbb R^{k+1}\) 中令
\[
 E=\bigcup_{q\in\mathbb Q}\bigl(\mathbb R^k\times\{q\}\bigr).
\]
证明 \(E\) 是可数 \(k\)-可整流 Borel 集，\(\mathcal H^k|_E\) 为 \(\sigma\) 有限测度，却在每个非空开球上取无穷值。
进一步证明：对任意 \(x\in E\)、任意 \(k\) 维线性子空间 \(T\subset\mathbb R^{k+1}\) 及 \(r>0\)，有
\[
 \mathcal H^k\Bigl(
 \bigl\{\,y\in E\cap B(x,r)\mid
 \operatorname{dist}(y-x,T)>\tfrac12|y-x|\,\bigr\}
 \Bigr)=\infty.
\]
据此说明该集合没有近似切平面，且不能删去密度与切平面定理中的局部有限性假设。
\end{exercise}
% 来源：自拟；用锥外球中的无穷个正面积切片实际检验近似切条件，而不从缺少假设直接推断结论失败。
\begin{solution}
每个平面 \(\mathbb R^k\times\{q\}\) 都是等距线性参数化的平移像，所以 \(E\) 可数 \(k\)-可整流。各平面闭、并为可数，故 \(E\) 是 Borel 集。再将每个平面分成与 \([-j,j]^{k+1}\) 的交，就得到可数个有限 \(\mathcal H^k\) 测度集的覆盖，证明了 \(\sigma\) 有限性。

任取开球 \(B(a,r)\)。对所有满足 \(|q-a_{k+1}|<r/2\) 的有理数 \(q\)，球与第 \(q\) 个平面的交是半径至少为 \(\sqrt3r/2\) 的 \(k\) 维球。这样的有理数无穷多，各切片互不相交，且每片的 \(\mathcal H^k\) 测度至少为 \(\omega_k(\sqrt3r/2)^k>0\)。可数可加性给出 \(\mathcal H^k\bigl(E\cap B(a,r)\bigr)=\infty\)。

现在固定 \(x,T,r\)，取垂直于 \(T\) 的单位向量 \(v\)，并令
\[
 U=B\left(x+\frac r2v,\frac r8\right).
\]
若 \(y\in U\)，则
\[
 |y-x|<\frac{5r}{8}<r,
 \quad \operatorname{dist}(y-x,T)>\frac{3r}{8}
 >\frac12|y-x|.
\]
所以 \(E\cap U\) 包含在题述锥外集合中。上一段已证明 \(\mathcal H^k(E\cap U)=\infty\)，所求断言随即成立。

近似切平面的锥外面积除以 \(r^k\) 应趋于零；这里对每个候选平面、每个半径，锥外面积已经是无穷。故近似切平面确实不存在。可数可整流性给出了参数片，却不保证这些片在小球中的总面积有限。
\end{solution}

\begin{exercise}\label{b1:ex:19-saddle-tangent-projection}
令
\[
 G(x,y)=(x,y,xy),
 \quad E=G([-1,1]^2),
 \quad P(x,y,z)=(y,z).
\]
这里 \(P\) 是投向第二、第三坐标平面的正交投影，随后将该平面等距识别为 \(\mathbb R^2\)。
\begin{enumerate}
\item\label{b1:item:19-saddle-tangent-normal}
求图像内部点 \(G(x,y)\) 的切平面及一个单位法向量，计算 \(J_2G\) 和切向面积因子 \(J_EP\)。
\item\label{b1:item:19-saddle-critical-fiber}
描述 \(P(E)\) 及投影的各纤维，指出切向面积因子为零的集合。用可整流集上的面积公式计算
\[
 \int_EJ_EP\,\mathrm d\mathcal H^2.
\]
\end{enumerate}
\end{exercise}
% 来源：自拟；通过非底面投影计算图像切向伸缩，保留被压成一点的整条临界纤维。
\begin{solution}
图像参数化的两列为
\[
 G_x=(1,0,y),
 \quad G_y=(0,1,x).
\]
它们线性无关，所以图像内部处处有切平面
\[
 T_{G(x,y)}E=\operatorname{span}\{(1,0,y),(0,1,x)\},
\]
并可选择单位法向量
\[
 n(x,y)=\frac{(-y,-x,1)}{\sqrt{1+x^2+y^2}}.
\]
Gram 行列式给出 \(J_2G=\sqrt{1+x^2+y^2}\)。另一方面，
\[
 (P\circ G)(x,y)=(y,xy),
 \quad D(P\circ G)=\begin{pmatrix}0&1\\y&x\end{pmatrix}.
\]
因此 \(J_2(P\circ G)=|y|\)。将投影限制在图像切平面上，再用线性面积因子的复合公式，得到
\[
 J_EP\bigl(G(x,y)\bigr)=
 \frac{|y|}{\sqrt{1+x^2+y^2}}.
\]
它也等于单位法向量第一分量的绝对值。图像边界是四条 Lipschitz 曲线的并，二维测度为零，不影响后续积分。

若把目标点记为 \((a,b)\)，原像条件是 \(y=a\)、\(xy=b\)，所以
\[
 P(E)=\{\,(a,b)\mid |a|\le1,\ |b|\le|a|\,\}.
\]
当 \(a\ne0\) 时，原像恰为 \(G(b/a,a)\)，是唯一一点；当 \((a,b)=(0,0)\) 时，整条线段 \(\{(x,0,0):-1\le x\le1\}\) 都在纤维中。没有其他非空纤维。

这条线段正是切向因子为零的集合，它的 \(\mathcal H^2\) 测度为零。除目标原点这一零测例外外，像内重数为 \(1\)，所以\cref{b1:thm:rectifiable-area-formula}给出
\[
 \int_EJ_EP\,\mathrm d\mathcal H^2
 =m_2\bigl(P(E)\bigr)=\int_{-1}^12|a|\,\mathrm da=2.
\]
也可由图像面积公式直接核对左端为 \(\int_{[-1,1]^2}|y|\,\mathrm dx\,\mathrm dy=2\)。投影存在无穷重数的纤维，并不妨碍几乎处处重数为 \(1\) 的面积计算。
\end{solution}

\begin{exercise}\label{b1:ex:19-sphere-projection-density}
设 \(S^2=\{\,(x,y,z)\in\mathbb R^3\mid x^2+y^2+z^2=1\,\}\)，令 \(P(x,y,z)=(x,y)\)。求切向因子 \(J_{S^2}P\)，验证
\[
 \int_{S^2}J_{S^2}P\,\mathrm d\mathcal H^2=2\pi.
\]
再求 \(P_\#(\mathcal H^2|_{S^2})\) 关于平面面积的密度，并说明单位圆周是否承受额外的推前质量。
\end{exercise}
% 来源：自拟；球面积在前文已知，此处训练切向 Jacobi、二重投影与推前密度，不重复球面积的参数化证明。
\begin{solution}
在点 \(p=(x,y,z)\in S^2\)，切平面为 \(p^\perp\)。取其定向标准正交基 \(v,w\)，使 \(v\times w=p\)。投影限制在该平面上的矩阵为
\[
 \begin{pmatrix}v_1&w_1\\v_2&w_2\end{pmatrix},
\]
其行列式是 \(v_1w_2-v_2w_1=(v\times w)_3=z\)。因此
\[
 J_{S^2}P(p)=|z|.
\]
对开单位圆盘中的每一点 \((a,b)\)，两个原像为
\[
 \left(a,b,\sqrt{1-a^2-b^2}\right),
 \quad \left(a,b,-\sqrt{1-a^2-b^2}\right).
\]
圆周上原像只有一个，圆盘外没有原像。圆周的二维测度为零，故切向面积公式给出
\[
 \int_{S^2}|z|\,\mathrm d\mathcal H^2
 =\int_{a^2+b^2<1}2\,\mathrm da\,\mathrm db=2\pi.
\]

赤道是 Lipschitz 曲线，\(\mathcal H^2\) 测度为零。在其补集上把 \(h\circ P/|z|\) 作为非负权重代入面积公式，就得到
\[
 \int_{S^2}h\bigl(P(p)\bigr)\,\mathrm d\mathcal H^2(p)
 =\int_{a^2+b^2<1}
 h(a,b)\frac{2}{\sqrt{1-a^2-b^2}}\,\mathrm da\,\mathrm db.
\]
所以密度为
\[
 \rho(a,b)=\frac{2}{\sqrt{1-a^2-b^2}},
 \quad a^2+b^2<1,
\]
在其余点另定义为零。极坐标给出
\[
 \int\rho\,\mathrm dm_2
 =4\pi\int_0^1\frac r{\sqrt{1-r^2}}\,\mathrm dr=4\pi,
\]
与球面积一致。密度趋近圆周时发散，但圆周的原像恰为赤道，原测度为零，所以不存在附加的圆周质量。
\end{solution}

\begin{exercise}\label{b1:ex:19-weighted-cross-and-atom}
在 \(\mathbb R^2\) 中令 \(L_1=\mathbb R\times\{0\}\)、\(L_2=\{0\}\times\mathbb R\)，并定义
\[
 E=L_1\cup L_2,
 \quad \mu=2\mathcal H^1|_{L_1}+3\mathcal H^1|_{L_2}.
\]
计算每个 \(x\in E\) 处的极限
\[
 \lim_{r\downarrow0}\frac{\mu\bigl(B(x,r)\bigr)}{2r},
\]
并与集合 \(E\) 在原点的密度比较。说明 \(\mu\) 为何是 \(1\)-可整流测度，而对任意 \(c>0\)，测度 \(\mu+c\delta_0\) 都不是 \(1\)-可整流测度，尽管二者有相同的支撑。
\end{exercise}
% 来源：自拟；显式区分集合密度、加权测度密度和仅集中在可整流集合上的测度。
\begin{solution}
在 \(L_1\setminus\{0\}\) 的点处，足够小的球不与 \(L_2\) 相交，且在 \(L_1\) 上截出长度 \(2r\) 的线段，所以 \(\mu\bigl(B(x,r)\bigr)=4r\)，所求极限为 \(2\)。在 \(L_2\setminus\{0\}\) 上同理得到极限 \(3\)。在原点，两条直线都贡献长度 \(2r\)，故
\[
 \mu\bigl(B(0,r)\bigr)=10r,
 \quad \lim_{r\downarrow0}\frac{\mu\bigl(B(0,r)\bigr)}{2r}=5.
\]
若去掉权重，\(\mathcal H^1\bigl(E\cap B(0,r)\bigr)=4r\)，所以集合密度是 \(2\)，不是测度密度 \(5\)。

集合 \(E\) 可数 \(1\)-可整流，且 \(\mathcal H^1|_E\) 局部有限。令 \(\theta=2\) 于 \(L_1\setminus\{0\}\)，\(\theta=3\) 于 \(L_2\setminus\{0\}\)，在原点任取有限值，便有
\[
 \mu=\theta\mathcal H^1|_E.
\]
因此 \(\mu\) 是 \(1\)-可整流 Radon 测度。原点的权值怎样修改都不改变 \(\mu\)，却不能改变原点处计算出的密度 \(5\)；密度等于权值的结论只要求几乎处处成立。

另一方面，\(\mathcal H^1(\{0\})=0\)，而 \((\mu+c\delta_0)(\{0\})=c>0\)。任何 \(\theta\mathcal H^1|_R\) 形式的测度都不可能在这个单点上赋正质量，故 \(\mu+c\delta_0\) 不是 \(1\)-可整流测度。它在原点的上述密度为无穷。两种测度的支撑都是 \(E\)，可见支撑是可整流集合并不足以代替关于 \(\mathcal H^1\) 的绝对连续性。
\end{solution}

\begin{exercise}\label{b1:ex:19-null-addition-topology}
令 \(S^1\subset\mathbb R^2\) 为单位圆周，并令 \(E=S^1\cup\mathbb Q^2\)。证明：
\begin{enumerate}
\item\label{b1:item:19-null-addition-measure}
集合 \(E\) 可数 \(1\)-可整流，且 \(\mathcal H^1|_E=\mathcal H^1|_{S^1}\)。两集合的近似切线在 \(\mathcal H^1\) 几乎处处意义下相同。
\item\label{b1:item:19-null-addition-closure}
有 \(\overline E=\mathbb R^2\)，但测度 \(\mathcal H^1|_E\) 的支撑是 \(S^1\)，并且 \(\overline E\) 不是可数 \(1\)-可整流集。
\end{enumerate}
\end{exercise}
% 来源：自拟；通过零测稠密附加集区分模零几何、拓扑闭包与测度支撑。
\begin{solution}
有理点集可数，所以 \(\mathcal H^1(\mathbb Q^2)=0\)。圆周是区间的 Lipschitz 像，故 \(E\) 可数 \(1\)-可整流。对任意 Borel 集 \(A\)，附加的有理点不改变一维测度，因此
\[
 \mathcal H^1(A\cap E)=\mathcal H^1(A\cap S^1).
\]
近似切线定义所用的小球质量和锥外质量也逐一相同。圆周的切线在 \(p\in S^1\) 处为 \(p^\perp\)，所以 \(E\) 与圆周具有相同的几乎处处近似切线。

由于 \(\mathbb Q^2\) 稠密，\(\overline E=\mathbb R^2\)。但是圆周外每个点都有一个不交圆周的小球，其 \(\mathcal H^1|_E\) 测度为零；圆周上每个点的任意小球都含正长度圆弧。于是测度的支撑恰为 \(S^1\)。

最后，每个 Lipschitz 曲线像的二维 Lebesgue 测度为零：将参数域分成有界部分，再用 Lipschitz 映射不增加 Hausdorff 维数即可。任意 \(\mathcal H^1\) 零测集也为 \(\mathcal H^2=m_2\) 零测集。因此每个可数 \(1\)-可整流集都为二维零测集，而 \(\mathbb R^2\) 不是。可整流性及其几乎处处切线不随零测修改而变化，但拓扑闭包可以彻底改变。
\end{solution}

\begin{exercise}\label{b1:ex:19-rectifiable-cylinder}
设 \(1\le k\le m\)、\(l\ge1\) 为整数。先证明：若 \(N\subset\mathbb R^m\) 满足 \(\mathcal H^k(N)=0\)，则
\[
 \mathcal H^{k+l}\bigl(N\times[0,1]^l\bigr)=0.
\]
再由此证明：若 \(E\subset\mathbb R^m\) 是可数 \(k\)-可整流 Borel 集，则 \(E\times[0,1]^l\) 是可数 \((k+l)\)-可整流集。
提示：不能把定义中未被参数片覆盖的零测余项直接省掉，须先处理它与立方体的乘积。
\end{exercise}
% 来源：自拟；用同尺度网格控制零测余项的柱体，再构造乘积参数片；不使用一般集合的错误维数相加上界。
\begin{solution}
给定 \(\delta\in(0,1)\) 和 \(\varepsilon>0\)。由 \(\mathcal H^k(N)=0\)，可用可数个半径 \(r_i\in(0,\delta)\) 的球 \(B_i\) 覆盖 \(N\)，并使
\[
 \sum_i r_i^k<\varepsilon.
\]
这可由任意小直径的 Hausdorff 覆盖分别套入球得到，归一化常数吸收到 \(\varepsilon\) 中。

对每个 \(i\)，把 \([0,1]^l\) 分成边长至多 \(r_i\) 的小立方体，所需数目不超过 \(2^lr_i^{-l}\)。每个乘积 \(B_i\times Q\) 的直径至多为 \(\sqrt{4+l}\,r_i\)。这些乘积覆盖 \(N\times[0,1]^l\)，其 \((k+l)\) 次直径和至多为
\[
 \sum_i2^lr_i^{-l}
 \bigl(\sqrt{4+l}\,r_i\bigr)^{k+l}
 \le C_{k,l}\sum_i r_i^k
 <C_{k,l}\varepsilon.
\]
令 \(\varepsilon\downarrow0\)，再令 \(\delta\downarrow0\)，即得第一问。

现按可数可整流性的定义，选 Lipschitz 映射 \(f_j:A_j\subset\mathbb R^k\rightarrow\mathbb R^m\)，使
\[
 N=E\setminus\bigcup_jf_j(A_j),
 \quad \mathcal H^k(N)=0.
\]
对每个 \(j\) 定义
\[
 F_j:A_j\times[0,1]^l\longrightarrow\mathbb R^{m+l},
 \quad F_j(a,b)=(f_j(a),b).
\]
映射 \(F_j\) 的 Lipschitz 常数不超过 \(\max\{\operatorname{Lip}(f_j),1\}\)。这些 \((k+l)\) 维参数片覆盖 \((E\setminus N)\times[0,1]^l\)，而剩余的 \(N\times[0,1]^l\) 已由第一问证明为零测集。这就建立了柱体的可数可整流性。证明中立方体的网格数 \(O(r_i^{-l})\) 抵消了额外的 \(l\) 次尺度因子。
\end{solution}

\begin{exercise}\label{b1:ex:19-pure-bilipschitz-decomposition}
设 \(1\le k\le n\)。证明纯 \(k\)-不可整流 Borel 集的 Borel 子集与可数并仍为纯 \(k\)-不可整流集。再设 \(F:\mathbb R^n\rightarrow\mathbb R^n\) 为双 Lipschitz 同胚，证明集合 \(P\) 纯 \(k\)-不可整流当且仅当 \(F(P)\) 纯 \(k\)-不可整流。

若 Borel 集 \(E\) 的 \(\mathcal H^k\) 测度 \(\sigma\) 有限，且 \(E=R\mathbin{\dot\cup}P\) 是本章的可整流—纯不可整流分解，证明 \(F(E)=F(R)\mathbin{\dot\cup}F(P)\) 给出 \(F(E)\) 的对应分解，其唯一性仍按 \(\mathcal H^k\) 零测集理解。
\end{exercise}
% 来源：自拟；把纯不可整流性的测试定义用于双 Lipschitz 变换，进而检验结构分解对坐标选择的独立性。
\begin{solution}
若 \(P'\subset P\)，则 \(P'\) 与任意 Lipschitz \(k\) 维像的交包含在 \(P\) 与该像的交中；后者零测便推出前者零测。若每个 \(P_j\) 纯 \(k\)-不可整流，则对任意 Lipschitz 映射 \(g:A\subset\mathbb R^k\rightarrow\mathbb R^n\)，
\[
 \mathcal H^k\Bigl(\bigcup_j\bigl(P_j\cap g(A)\bigr)\Bigr)
 \le\sum_j\mathcal H^k\bigl(P_j\cap g(A)\bigr)=0.
\]
这证明子集与可数并的封闭性。

假设 \(P\) 纯 \(k\)-不可整流。由于 \(F^{-1}\circ g\) 仍为 Lipschitz 映射，集合
\[
 P\cap(F^{-1}\circ g)(A)
\]
为 \(\mathcal H^k\) 零测集。又有
\[
 F(P)\cap g(A)=F\bigl(P\cap(F^{-1}\circ g)(A)\bigr).
\]
Lipschitz 映射把 \(\mathcal H^k\) 零测集送到零测集，所以右端为零测集，证明 \(F(P)\) 纯 \(k\)-不可整流。对 \(F^{-1}\) 重复此论证便得逆向蕴涵。

同时，若 \(R\) 可数 \(k\)-可整流，将它的各参数映射与 \(F\) 复合即可覆盖 \(F(R)\)，原来的零测余项仍为零测。因此 \(F(R)\) 可数 \(k\)-可整流。\(F\) 单射保证两部分不交，且它们覆盖 \(F(E)\)。若 \(E\) 被可数个有限 \(\mathcal H^k\) 测度集覆盖，Lipschitz 测度估计也保证 \(F(E)\) 仍为 \(\sigma\) 有限。故本章结构分解的唯一性适用，任何另一分解都只在 \(\mathcal H^k\) 零测集上与 \(F(R),F(P)\) 不同。

这里同时使用了两个方向的 Lipschitz 控制：前向控制保持零测集，反向控制把目标中的每个测试参数片拉回为合法的测试参数片。
\end{solution}

\begin{exercise}\label{b1:ex:19-rectifiable-weak-limit-dimension}
对整数 \(N\ge1\)，令
\[
 S_{j,N}=[0,1]\times\left\{\frac{j-1/2}{N}\right\},
 \quad j=1,\ldots,N,
 \quad \mu_N=\frac1N\sum_{j=1}^N\mathcal H^1|_{S_{j,N}}.
\]
\begin{enumerate}
\item\label{b1:item:19-grid-weak-limit}
证明 \(\mu_N\) 是 \(1\)-可整流概率测度，并且弱收敛到 \(\mu=m_2|_{[0,1]^2}\)。证明极限 \(\mu\) 不是 \(1\)-可整流测度。
\item\label{b1:item:19-grid-density-bound}
求 \(\mu_N\) 在其 \(\mathcal H^1\) 几乎每个支撑点处的一维密度，并证明对每个 \(a\in\mathbb R^2\)、\(r>0\)，
\[
 \mu_N\bigl(B(a,r)\bigr)\le4r^2+\frac{4r}{N}.
\]
据此说明：共同紧支撑、固定总质量，甚至统一的一维球增长上界，都不足以保证弱极限仍为 \(1\)-可整流测度。
\end{enumerate}
\end{exercise}
% 来源：自拟；以加权网格线的 Riemann 和展示可整流测度的弱极限可增加维数，不把 Preiss 概览当作证明工具。
\begin{solution}
每条线段长度为 \(1\)，彼此不交，所以 \(\mu_N\) 的总质量为 \(1\)。它等于有限可整流集 \(\bigcup_jS_{j,N}\) 上的常权 \((1/N)\mathcal H^1\)，故为 \(1\)-可整流概率测度。

设 \(h\in C_b(\mathbb R^2)\)。它在单位方形上一致连续，记其在该方形上的连续模为
\[
 \omega_h(\delta)=\sup\{\,|h(x)-h(y)|\mid x,y\in[0,1]^2,\ |x-y|\le\delta\,\}.
\]
把方形按第二坐标分成 \(N\) 条等宽条带，每条带中的点到其中线的距离不超过 \(1/(2N)\)。于是
\[
 \left|\int h\,\mathrm d\mu_N-\int_{[0,1]^2}h\,\mathrm dm_2\right|
 \le\omega_h\left(\frac1{2N}\right)\longrightarrow0.
\]
这正是 \(\mu_N\Rightarrow\mu\)。

每个可数 \(1\)-可整流集都是二维零测集，理由见\cref{b1:ex:19-null-addition-topology}。如果 \(\mu\) 是 \(1\)-可整流测度，它必须集中在某个这样的集合 \(R\) 上；但 \(\mu(R)=m_2(R\cap[0,1]^2)=0\)，与总质量 \(1\) 矛盾。因此弱极限不再 \(1\)-可整流。

对固定 \(N\)，在任一线段的非端点处，半径充分小的球只与该线段相交，并截出长度 \(2r\) 的区间。故
\[
 \lim_{r\downarrow0}\frac{\mu_N\bigl(B(x,r)\bigr)}{2r}=\frac1N
\]
在 \(\mu_N\) 几乎处处成立。另一方面，与 \(B(a,r)\) 相交的水平线条数不超过 \(2Nr+2\)，每条在球中留下的长度不超过 \(2r\)，所以
\[
 \mu_N\bigl(B(a,r)\bigr)\le\frac{2r}{N}(2Nr+2)
 =4r^2+\frac{4r}{N}.
\]
当 \(0<r\le1\) 时，右端不超过 \(8r\)；当 \(r>1\) 时，由总质量 \(1\) 同样得到 \(\mu_N\bigl(B(a,r)\bigr)\le8r\)。因此存在独立于 \(N\) 的一维球增长上界。

但各条线上的正密度为 \(1/N\)，没有统一的正下界；极限测度在方形内部满足 \(\mu\bigl(B(x,r)\bigr)/(2r)=\pi r/2\to0\)。球增长的上界容许质量铺展到更高维区域，不能保证极限的一维密度仍为正。
\end{solution}
